% ============================================================
%  module_02_cell_biochem.tex  —  IBO Biology Platform
%  Module 2: Cell Biology and Biochemistry
%
%  STANDALONE: compile this file directly on Overleaf.
%  Upload together with: preamble.tex + the four .jpg images.
%  Compiler: pdfLaTeX. Root: module_02_cell_biochem.tex
% ============================================================

\documentclass[12pt, a4paper, oneside]{report}
% ============================================================
%  preamble.tex  —  IBO Biology Platform
%  Minimal academic style. pdfLaTeX compatible.
% ============================================================

% --- Page geometry ---
\usepackage[a4paper,
            top=2.5cm, bottom=2.8cm,
            left=3cm, right=2.5cm]{geometry}

% --- Typography ---
\usepackage[T1]{fontenc}
\usepackage[utf8]{inputenc}
\usepackage{lmodern}
\usepackage{microtype}
\usepackage{setspace}
\setstretch{1.12}

% --- Mathematics ---
\usepackage{amsmath}
\usepackage{amssymb}
\usepackage{amsfonts}
\usepackage{mathtools}
\usepackage{bm}
\usepackage{siunitx}
\sisetup{per-mode=symbol}

% --- Chemistry ---
\usepackage[version=4]{mhchem}

% --- Graphics ---
\usepackage{graphicx}
\usepackage{xcolor}
\usepackage{tikz}
\usetikzlibrary{calc}

% --- Floats and captions ---
\usepackage{float}
\usepackage{caption}
\captionsetup{font=small, labelfont=bf, labelsep=period, width=0.92\textwidth}

% --- Tables ---
\usepackage{booktabs}
\usepackage{tabularx}
\usepackage{array}

% --- Lists ---
\usepackage{enumitem}
\setlist[enumerate,1]{label=\textbf{(\alph*)}, itemsep=2pt}
\setlist[enumerate,2]{label=\roman*.,  itemsep=1pt}
\setlist[itemize]{itemsep=2pt}

% --- Framed environments (for solutions, notes) ---
\usepackage{mdframed}

% --- Colour palette: one accent, otherwise near-black ---
\definecolor{accent}{HTML}{1A2E4A}    % dark navy
\definecolor{probback}{HTML}{F5F6F7}  % near-white box background
\definecolor{stargold}{HTML}{A07820}  % muted gold
\definecolor{rulecol}{HTML}{AAAAAA}   % light rule colour
\definecolor{codebg}{HTML}{F8F8F8}
\definecolor{codeframe}{HTML}{CCCCCC}
\definecolor{codekw}{HTML}{00397A}
\definecolor{codestr}{HTML}{B83200}
\definecolor{codecomment}{HTML}{4A7A4A}
\definecolor{codenumber}{HTML}{999999}

% --- Hyperlinks ---
\usepackage{hyperref}
\hypersetup{
    colorlinks = true,
    linkcolor  = accent,
    citecolor  = accent,
    urlcolor   = accent,
    pdftitle   = {IBO Biology Platform},
}

% --- Code listings ---
\usepackage{listings}
\usepackage{lstautogobble}

\lstdefinestyle{python}{
    language          = Python,
    backgroundcolor   = \color{codebg},
    basicstyle        = \ttfamily\small,
    keywordstyle      = \color{codekw}\bfseries,
    stringstyle       = \color{codestr},
    commentstyle      = \color{codecomment}\itshape,
    numberstyle       = \tiny\color{codenumber},
    numbers           = left,
    stepnumber        = 1,
    numbersep         = 8pt,
    frame             = single,
    framesep          = 4pt,
    rulecolor         = \color{codeframe},
    breaklines        = true,
    showstringspaces  = false,
    tabsize           = 4,
    captionpos        = b,
    autogobble        = true,
    morekeywords      = {as, with, None, True, False,
                         np, plt, odeint, linspace, exp, log, pi, array},
}
\lstset{style=python}

% --- tcolorbox (problems only) ---
\usepackage{tcolorbox}
\tcbuselibrary{breakable, skins}

% Problem box: clean, numbered, minimal
\newtcolorbox[auto counter, number within=section]{problem}[1]{%
    breakable,
    enhanced,
    colback    = probback,
    colframe   = accent,
    boxrule    = 0.5pt,
    arc        = 1pt,
    fonttitle  = \bfseries\sffamily\small,
    coltitle   = white,
    colbacktitle = accent,
    title      = {Problem~\thetcbcounter.\enspace #1},
    attach boxed title to top left={yshift=-1.5mm, xshift=3mm},
    boxed title style={arc=1pt, boxrule=0pt},
    top        = 5pt, bottom = 5pt,
    left       = 6pt, right  = 6pt,
}

% Solution: simple left-rule indent, no background
\newenvironment{solution}{%
    \begin{mdframed}[
        linewidth    = 1.5pt,
        linecolor    = accent!40,
        topline      = false,
        rightline    = false,
        bottomline   = false,
        leftmargin   = 4pt,
        innerleftmargin = 10pt,
        innertopmargin  = 4pt,
        innerbottommargin = 4pt,
        backgroundcolor = white,
    ]%
    \noindent\textit{Solution.}\enspace\ignorespaces
}{%
    \end{mdframed}%
}

% Historical note: small italic, thin left rule
\newenvironment{historicalnote}[1]{%
    \begin{mdframed}[
        linewidth    = 2pt,
        linecolor    = rulecol,
        topline      = false,
        rightline    = false,
        bottomline   = false,
        leftmargin   = 4pt,
        innerleftmargin = 10pt,
        innertopmargin  = 3pt,
        innerbottommargin = 3pt,
        backgroundcolor = white,
    ]%
    {\small\sffamily\bfseries Historical note: #1}\\[3pt]%
    \small\itshape\ignorespaces
}{%
    \end{mdframed}%
}

% Definition: bold label, thin left rule
\newenvironment{definition}[1]{%
    \begin{mdframed}[
        linewidth    = 1.5pt,
        linecolor    = accent!60,
        topline      = false,
        rightline    = false,
        bottomline   = false,
        leftmargin   = 4pt,
        innerleftmargin = 10pt,
        innertopmargin  = 3pt,
        innerbottommargin = 3pt,
        backgroundcolor = white,
    ]%
    \noindent\textbf{Definition (#1).}\enspace\ignorespaces
}{%
    \end{mdframed}%
}

% Key remark: just bold inline label, no box
\newenvironment{keyinsight}{%
    \par\medskip\noindent\textbf{Remark.}\enspace\ignorespaces
}{\par\medskip}

% Model limitations: just bold inline label
\newenvironment{limitation}{%
    \par\medskip\noindent\textbf{Limitations of this model.}\enspace\ignorespaces
}{\par\medskip}

% Summary: simple ruled block
\newenvironment{summary}{%
    \begin{mdframed}[
        linewidth = 0.5pt, linecolor = accent!50,
        leftmargin = 0pt, innerleftmargin = 8pt,
        innertopmargin = 6pt, innerbottommargin = 6pt,
        backgroundcolor = probback,
    ]%
    \noindent\textbf{Summary.}\enspace\ignorespaces
}{%
    \end{mdframed}%
}

% Code lab: just a listing caption prefix
\newenvironment{codelab}[1]{%
    \par\medskip\noindent\textbf{Code.}\enspace\textit{#1}\par\smallskip
}{\par}

% ============================================================
%  Difficulty stars — plain \ifnum, no loops
% ============================================================
\usepackage{pifont}

\newcommand{\starfilled}{{\color{stargold}\ding{72}}}
\newcommand{\starempty} {{\color{stargold!40}\ding{73}}}

\newcommand{\stars}[1]{%
  \ifnum#1>0\starfilled\else\starempty\fi
  \ifnum#1>1\starfilled\else\starempty\fi
  \ifnum#1>2\starfilled\else\starempty\fi
  \ifnum#1>3\starfilled\else\starempty\fi
  \ifnum#1>4\starfilled\else\starempty\fi
}

% ============================================================
%  Discipline tags — plain \fcolorbox, safe everywhere
% ============================================================
\newcommand{\disciplinetag}[1]{%
  {\setlength{\fboxsep}{2pt}\setlength{\fboxrule}{0.4pt}%
   \fcolorbox{accent!50}{accent!8}{%
     \textcolor{accent}{\sffamily\tiny\bfseries\vphantom{Ag}#1}}}%
}

% ============================================================
%  Section / chapter headings — clean, minimal
% ============================================================
\usepackage{titlesec}

\titleformat{\chapter}[block]
  {\normalfont\LARGE\bfseries\sffamily\color{accent}}
  {Module~\thechapter.}{0.6em}{}[\vspace{2pt}\rule{\textwidth}{0.4pt}]

\titleformat{\section}
  {\normalfont\large\bfseries\sffamily\color{accent}}
  {\thesection}{0.7em}{}

\titleformat{\subsection}
  {\normalfont\normalsize\bfseries\sffamily}
  {\thesubsection}{0.6em}{}

\titleformat{\subsubsection}
  {\normalfont\normalsize\itshape}
  {}{0em}{}

% ============================================================
%  Header / footer — minimal
% ============================================================
\usepackage{fancyhdr}
\pagestyle{fancy}
\fancyhf{}
\fancyhead[RO,LE]{\small\sffamily\thepage}
\fancyhead[LO]{\small\sffamily\nouppercase{\rightmark}}
\fancyhead[RE]{\small\sffamily\nouppercase{\leftmark}}
\renewcommand{\headrulewidth}{0.3pt}
\renewcommand{\footrulewidth}{0pt}

% ============================================================
%  Maths macros
% ============================================================
\newcommand{\kB}{k_{\mathrm{B}}}
\newcommand{\NA}{N_{\mathrm{A}}}
\newcommand{\dG}{\Delta G}
\newcommand{\ddG}{\Delta\Delta G}
\newcommand{\dH}{\Delta H}
\newcommand{\dS}{\Delta S}
\newcommand{\Kd}{K_{\mathrm{d}}}
\newcommand{\Ka}{K_{\mathrm{a}}}
\newcommand{\Keq}{K_{\mathrm{eq}}}
\newcommand{\vmax}{V_{\mathrm{max}}}
\newcommand{\Km}{K_{\mathrm{m}}}
\newcommand{\conc}[1]{[\mathrm{#1}]}
\newcommand{\od}[2]{\frac{\mathrm{d}#1}{\mathrm{d}#2}}
\newcommand{\pd}[2]{\frac{\partial #1}{\partial #2}}
\newcommand{\e}{\mathrm{e}}
\newcommand{\abs}[1]{\left|#1\right|}

% ============================================================
%  Text abbreviations
% ============================================================
\newcommand{\Ecoli}{\textit{E.~coli}}
\newcommand{\invitro}{\textit{in vitro}}
\newcommand{\invivo}{\textit{in vivo}}
\newcommand{\etal}{\textit{et al.}}

% End of preamble.tex


\begin{document}

% ---------- Title block ----------
\begin{titlepage}
  \centering
  \vspace*{4cm}
  {\sffamily\fontsize{13}{16}\selectfont\color{accent} IBO Biology Platform}\\[0.6em]
  {\rule{0.4\textwidth}{0.4pt}}\\[0.6em]
  {\sffamily\fontsize{22}{28}\selectfont\bfseries Module 2}\\[0.4em]
  {\sffamily\fontsize{16}{20}\selectfont\bfseries\color{accent}
    Cell Biology and Biochemistry}\\[1.2em]
  {\small\color{gray}
    Enzyme kinetics \textbullet\ Inhibition \textbullet\
    Bioenergetics \textbullet\ Oxidative phosphorylation \textbullet\
    Metabolic control \textbullet\ Signal transduction}\\[0.6em]
  {\rule{0.4\textwidth}{0.4pt}}\\[0.5em]
  {\small\color{gray} \today}
\end{titlepage}

\tableofcontents
\newpage

% ============================================================
%  Chapter begins here
% ============================================================

\chapter{Cell Biology and Biochemistry}
\label{ch:cell_biochem}

\begin{tcolorbox}[colback=probback, colframe=accent, boxrule=0.5pt, arc=1pt,
  fonttitle=\bfseries\sffamily\small, coltitle=white, colbacktitle=accent,
  title={Module Overview}, top=4pt, bottom=4pt, left=6pt, right=6pt]
  \textbf{Topics:} Enzyme kinetics and inhibition;
  bioenergetics and ATP hydrolysis free energy;
  the electron transport chain and chemiosmosis;
  F$_1$F$_0$ ATP synthase stoichiometry;
  metabolic flux control; signal transduction kinetics.\\[3pt]
  \textbf{Prerequisites:} Thermodynamics ($\dG$, equilibrium);
  basic kinetics (rate laws, steady state); Module~1.\\[3pt]
  \textbf{Problem set:} 10 problems, \stars{2}--\stars{5}.
\end{tcolorbox}

% ============================================================
\section{Enzyme kinetics}
% ============================================================

\subsection{The catalytic problem}

Enzymes accelerate reactions by factors of $10^6$--$10^{20}$ relative to
the uncatalysed rate. They do this without altering the equilibrium
position — they lower the activation energy $\Delta G^\ddagger$ by
stabilising the transition state. The central question is: how does
reaction rate depend on substrate concentration?

\begin{historicalnote}{Michaelis and Menten (1913); Briggs and Haldane (1925)}
  Leonor Michaelis and Maud Menten studied invertase (sucrose hydrolysis)
  in 1913, establishing the hyperbolic dependence of rate on substrate
  concentration and defining $K_M$ under the rapid-equilibrium assumption
  ($k_2 \ll k_{-1}$, so $K_M = K_d$). George Briggs and J.B.S. Haldane
  (1925) generalised the derivation using the steady-state assumption for
  [ES], removing the restriction on rate constants. The result is the
  same equation, but with a different physical meaning for $K_M$.
\end{historicalnote}

\subsection{Derivation of the Michaelis-Menten equation}

The minimal mechanism:
\begin{equation}
  \mathrm{E} + \mathrm{S}
    \underset{k_{-1}}{\overset{k_1}{\rightleftharpoons}}
  \mathrm{ES}
    \overset{k_\text{cat}}{\longrightarrow}
  \mathrm{E} + \mathrm{P}.
  \label{eq:mmechanism}
\end{equation}

\textbf{Assumptions.}
\begin{enumerate}
  \item $[\mathrm{S}] \gg [\mathrm{E}]_0$ (total enzyme concentration),
        so $[\mathrm{S}]$ is approximately constant (not depleted by ES
        formation).
  \item The reaction is measured at initial velocity $v_0$, before
        significant product accumulates, so the reverse reaction
        $\mathrm{P} \to \mathrm{ES}$ is negligible.
  \item The steady-state assumption: $\od{[\mathrm{ES}]}{t} \approx 0$.
\end{enumerate}

\textbf{Derivation.}
Rate of ES formation: $k_1[\mathrm{E}][\mathrm{S}]$.\\
Rate of ES disappearance: $(k_{-1} + k_\text{cat})[\mathrm{ES}]$.\\
Steady state:
\[
  k_1[\mathrm{E}][\mathrm{S}] = (k_{-1} + k_\text{cat})[\mathrm{ES}].
\]
Conservation of total enzyme: $[\mathrm{E}]_0 = [\mathrm{E}] + [\mathrm{ES}]$,
so $[\mathrm{E}] = [\mathrm{E}]_0 - [\mathrm{ES}]$. Substituting:
\[
  k_1([\mathrm{E}]_0 - [\mathrm{ES}])[\mathrm{S}]
    = (k_{-1} + k_\text{cat})[\mathrm{ES}].
\]
Define $K_M \equiv (k_{-1} + k_\text{cat})/k_1$. Solving for [ES]:
\[
  [\mathrm{ES}] = \frac{[\mathrm{E}]_0[\mathrm{S}]}{[\mathrm{S}] + K_M}.
\]
The reaction velocity $v = k_\text{cat}[\mathrm{ES}]$, and the maximum
velocity $V_{\max} = k_\text{cat}[\mathrm{E}]_0$ (all enzyme in ES form).
Therefore:

\begin{equation}
  \boxed{v = \frac{V_{\max}[\mathrm{S}]}{[\mathrm{S}] + K_M}}
  \label{eq:mm}
\end{equation}

\begin{figure}[H]
  \centering
  \includegraphics[width=0.92\textwidth]{michaelis_menten_kinetics_1786956238664.jpg}
  \caption{Left: the Michaelis-Menten reaction scheme and the
           hyperbolic $v$ vs.\ $[\mathrm{S}]$ curve. $K_M$ is
           the substrate concentration giving $v = V_{\max}/2$.
           Right: Lineweaver-Burk double-reciprocal transformation
           ($1/v$ vs.\ $1/[\mathrm{S}]$) linearises the data for
           determination of $V_{\max}$ and $K_M$ from intercepts.}
  \label{fig:mm}
\end{figure}

\subsection{Interpreting $K_M$ and $V_{\max}$}

At $[\mathrm{S}] = K_M$: $v = V_{\max}/2$ by inspection of
equation~\eqref{eq:mm}. So $K_M$ is operationally the substrate
concentration at half-maximal velocity. Mechanistically:
\[
  K_M = \frac{k_{-1} + k_\text{cat}}{k_1}.
\]
$K_M$ equals the true dissociation constant $K_d = k_{-1}/k_1$ only
when $k_\text{cat} \ll k_{-1}$ (rapid-equilibrium limit, the original
Michaelis assumption). In general, $K_M \geq K_d$.

$V_{\max} = k_\text{cat}[\mathrm{E}]_0$ depends on enzyme concentration,
so it is not an intrinsic property. $k_\text{cat}$ is intrinsic —
the first-order rate constant for ES breakdown, with units s$^{-1}$.
Typical $k_\text{cat}$ values range from $<1$~s$^{-1}$ (slow
isomerases) to $\sim10^7$~s$^{-1}$ (carbonic anhydrase).

\subsection{Catalytic efficiency and the diffusion limit}

The ratio $k_\text{cat}/K_M$ is the apparent second-order rate constant
for the reaction $\mathrm{E} + \mathrm{S} \to \mathrm{E} + \mathrm{P}$
at low substrate concentrations ($[\mathrm{S}] \ll K_M$):
\[
  v \approx \frac{k_\text{cat}}{K_M}[\mathrm{E}][\mathrm{S}]
  \quad \text{when } [\mathrm{S}] \ll K_M.
\]
$k_\text{cat}/K_M$ is bounded above by the diffusion-limited encounter
rate, $\sim 10^8$--$10^9$~M$^{-1}$s$^{-1}$. Enzymes at or near this
limit are called \textbf{catalytically perfect}: every encounter with
substrate leads to product. Examples include triosephosphate isomerase
($k_\text{cat}/K_M \approx 2\times10^8$~M$^{-1}$s$^{-1}$) and
acetylcholinesterase ($\approx 1.5\times10^8$~M$^{-1}$s$^{-1}$).

% ============================================================
\section{Enzyme inhibition}
% ============================================================

\subsection{Competitive inhibition}

A competitive inhibitor I binds only free enzyme E (not ES), at a site
that prevents substrate binding (often the active site):
\[
  \mathrm{E} + \mathrm{I} \underset{K_i}{\rightleftharpoons} \mathrm{EI},
  \qquad K_i = \frac{[\mathrm{E}][\mathrm{I}]}{[\mathrm{EI}]}.
\]
The enzyme is now partitioned between E, ES, and EI. The total enzyme
conservation becomes:
$[\mathrm{E}]_0 = [\mathrm{E}] + [\mathrm{ES}] + [\mathrm{EI}]$.
Using $[\mathrm{EI}] = [\mathrm{E}][\mathrm{I}]/K_i$, the steady-state
derivation gives:
\begin{equation}
  v = \frac{V_{\max}[\mathrm{S}]}{[\mathrm{S}] + K_M^\text{app}},
  \qquad
  K_M^\text{app} = K_M\!\left(1 + \frac{[\mathrm{I}]}{K_i}\right).
  \label{eq:competitive}
\end{equation}
$V_{\max}$ is unchanged (sufficient substrate outcompetes I). $K_M$ is
increased by the factor $(1 + [\mathrm{I}]/K_i)$. On a Lineweaver-Burk
plot: lines with different $[\mathrm{I}]$ intersect on the $y$-axis
(same $1/V_{\max}$, different slopes).

\subsection{Uncompetitive inhibition}

An uncompetitive inhibitor binds only ES (not free E):
\[
  \mathrm{ES} + \mathrm{I} \underset{K_i'}{\rightleftharpoons} \mathrm{ESI}.
\]
This removes ES from the productive pathway, pulling the E+S equilibrium
toward ES formation and paradoxically lowering the apparent $K_M$:
\begin{equation}
  V_{\max}^\text{app} = \frac{V_{\max}}{1 + [\mathrm{I}]/K_i'},
  \qquad
  K_M^\text{app} = \frac{K_M}{1 + [\mathrm{I}]/K_i'}.
  \label{eq:uncompetitive}
\end{equation}
Both $K_M^\text{app}$ and $V_{\max}^\text{app}$ decrease by the same
factor $\alpha' = 1 + [\mathrm{I}]/K_i'$. On a Lineweaver-Burk plot,
lines are parallel (same slope, different intercepts).

\subsection{Mixed inhibition}

A mixed inhibitor binds both E and ES, with distinct affinities $K_i$
and $K_i'$. The general result:
\begin{equation}
  v = \frac{V_{\max}[\mathrm{S}]}{\alpha K_M + \alpha'[\mathrm{S}]},
  \qquad
  \alpha = 1 + \frac{[\mathrm{I}]}{K_i},
  \quad
  \alpha' = 1 + \frac{[\mathrm{I}]}{K_i'}.
  \label{eq:mixed}
\end{equation}
When $\alpha = \alpha'$ (equal affinities for E and ES), this is
\textbf{noncompetitive} inhibition: $K_M$ is unchanged but $V_{\max}$
decreases. Lineweaver-Burk lines intersect on the $x$-axis.

% ============================================================
\section{Bioenergetics}
% ============================================================

\subsection{Free energy of ATP hydrolysis}

The standard free energy of ATP hydrolysis at pH 7, 25$^\circ$C,
$[\mathrm{Mg}^{2+}] = 1$~mM:
\[
  \mathrm{ATP} + \mathrm{H}_2\mathrm{O}
  \longrightarrow \mathrm{ADP} + \mathrm{P}_i,
  \qquad \dG^{\circ'} = -30.5~\text{kJ~mol}^{-1}.
\]
The actual free energy under cellular conditions:
\begin{equation}
  \dG = \dG^{\circ'} + RT\ln\frac{[\mathrm{ADP}][\mathrm{P}_i]}
        {[\mathrm{ATP}]}.
  \label{eq:dgatp}
\end{equation}
Typical mammalian cytoplasm: $[\mathrm{ATP}] \approx 5$~mM,
$[\mathrm{ADP}] \approx 0.5$~mM, $[\mathrm{P}_i] \approx 5$~mM.
At $T = 310$~K, $R = 8.314$~J~mol$^{-1}$~K$^{-1}$:
\[
  \dG = -30\,500 + 8.314 \times 310 \times
        \ln\!\left(\frac{0.5\times10^{-3} \times 5\times10^{-3}}
                        {5\times10^{-3}}\right).
\]
\[
  = -30\,500 + 2577 \times \ln(5\times10^{-4})
  = -30\,500 + 2577 \times (-7.60)
  \approx -50\text{ to }-54~\text{kJ~mol}^{-1}.
\]
This is $\approx 20$--22~$k_\mathrm{B}T$ per ATP hydrolysis — the actual
free energy available to drive cellular work.

\subsection{Coupled reactions}

A thermodynamically unfavourable reaction ($\dG_1 > 0$) can proceed if
coupled to ATP hydrolysis such that the overall $\dG_\text{total} < 0$:
\[
  \dG_\text{total} = \dG_1 + \dG_\text{ATP} < 0.
\]
The mechanism of coupling is never direct energy transfer — it is the
common sharing of a chemical intermediate. For example, in hexokinase,
glucose is phosphorylated using the $\gamma$-phosphate of ATP, producing
glucose-6-phosphate as a shared intermediate. The reaction:
\[
  \text{Glucose} + \text{ATP} \longrightarrow
  \text{Glucose-6-P} + \text{ADP},
  \quad \dG^{\circ'} = -16.7~\text{kJ~mol}^{-1}.
\]

\subsection{Phosphoryl transfer potential}

Free energies of hydrolysis for biologically common phosphorylated
compounds span a wide range. Selected values at pH~7, 25$^\circ$C:

\begin{center}
\begin{tabular}{lc}
\toprule
Compound & $\dG^{\circ'}$ (kJ~mol$^{-1}$) \\
\midrule
Phosphoenolpyruvate (PEP) & $-61.9$ \\
1,3-bisphosphoglycerate  & $-49.3$ \\
ATP                      & $-30.5$ \\
Glucose-1-phosphate      & $-20.9$ \\
Glucose-6-phosphate      & $-13.8$ \\
\bottomrule
\end{tabular}
\end{center}

ATP sits at an intermediate position, enabling it to accept phosphate
from high-energy donors (PEP, 1,3-BPG) and donate phosphate to
lower-energy acceptors (glucose, glycerol). This is the chemical basis
of ATP as the universal free energy currency.

% ============================================================
\section{Oxidative phosphorylation}
% ============================================================

\subsection{The electron transport chain}

\begin{figure}[H]
  \centering
  \includegraphics[width=0.95\textwidth]{electron_transport_chain_1786956256474.jpg}
  \caption{The mitochondrial electron transport chain. Complexes
           I--IV transfer electrons from NADH and FADH$_2$ to O$_2$,
           pumping protons into the intermembrane space. The resulting
           electrochemical gradient drives ATP synthase (Complex~V).
           Ubiquinone (Q) and cytochrome~$c$ (cyt~$c$) are mobile
           electron carriers.}
  \label{fig:etc}
\end{figure}

Electrons from NADH and FADH$_2$ pass through four protein complexes
embedded in the inner mitochondrial membrane:

\begin{enumerate}
  \item \textbf{Complex I} (NADH dehydrogenase, 44 subunits in humans):
        $\mathrm{NADH} \to \mathrm{NAD}^+ + \mathrm{H}^+ + 2e^-$.
        Electrons reduce ubiquinone (Q) to ubiquinol (QH$_2$).
        4~H$^+$ pumped per 2~electrons.
  \item \textbf{Complex II} (succinate dehydrogenase): oxidises
        succinate to fumarate, reducing Q to QH$_2$.
        \textbf{No proton pumping.}
  \item \textbf{Complex III} (cytochrome $bc_1$): QH$_2$ is oxidised,
        electrons pass via the Q-cycle to cytochrome~$c$.
        4~H$^+$ pumped per 2~electrons (Q-cycle mechanism).
  \item \textbf{Complex IV} (cytochrome $c$ oxidase): four electrons
        from cytochrome~$c$ reduce O$_2$ to 2~H$_2$O.
        4~H$^+$ pumped per 4~electrons (2~H$^+$ per pair).
\end{enumerate}

Proton stoichiometry per NADH: $4 + 4 + 2 = 10$~H$^+$.\\
Per FADH$_2$ (bypasses Complex I): $4 + 2 = 6$~H$^+$.

\begin{historicalnote}{Peter Mitchell and the chemiosmotic hypothesis (1961, Nobel 1978)}
  In 1961, Peter Mitchell proposed that the free energy of electron
  transfer is stored as an electrochemical proton gradient across the
  inner mitochondrial membrane, rather than as a direct chemical
  intermediate. This was highly controversial: the biochemical
  community expected a phosphorylated intermediate analogous to
  glycolysis. Mitchell's hypothesis was confirmed by the 1970s through
  reconstitution experiments showing that artificially imposed proton
  gradients drive ATP synthesis in isolated membrane vesicles.
  He received the Nobel Prize in Chemistry in 1978.
\end{historicalnote}

\subsection{The proton motive force}

The proton motive force (PMF), in units of volts, is:
\begin{equation}
  \Delta p = \Delta\psi - \frac{2.303\,RT}{F}\,\Delta\mathrm{pH},
  \label{eq:pmf}
\end{equation}
where $\Delta\psi$ is the electrical potential difference (matrix minus
intermembrane space, typically $-150$ to $-180$~mV in active
mitochondria) and $\Delta\mathrm{pH} = \mathrm{pH}_\text{matrix} -
\mathrm{pH}_\text{IMS} \approx +0.5$ to $+1.0$ units. At 310~K:
\[
  \frac{2.303\,RT}{F} = \frac{2.303 \times 8.314 \times 310}{96{,}485}
  = 61.5~\text{mV per pH unit.}
\]
With $\Delta\psi = -180$~mV and $\Delta\mathrm{pH} = 0.7$:
\[
  \Delta p = -180 - 61.5\times 0.7 = -180 - 43 = -223~\text{mV}.
\]
The free energy required to move one proton from the IMS to the matrix
against this gradient:
\begin{equation}
  \dG_\text{H}^+ = F\,\Delta p = 96{,}485 \times (-0.223)
  = -21.5~\text{kJ~mol}^{-1}.
  \label{eq:dgh}
\end{equation}

\subsection{F$_1$F$_0$ ATP synthase: rotary mechanism}

\begin{figure}[H]
  \centering
  \includegraphics[width=0.60\textwidth]{atp_synthase_f1f0_1786956248513.jpg}
  \caption{F$_1$F$_0$ ATP synthase. The F$_0$ domain (membrane-embedded
           $c$-ring and $a$-subunit) converts the proton flow into rotation
           of the $\gamma$ subunit. The F$_1$ domain ($\alpha_3\beta_3$
           hexamer) operates three catalytic $\beta$ sites in the
           conformational states Open~(O), Loose~(L), and Tight~(T),
           cycling sequentially as $\gamma$ rotates. Each 360$^\circ$
           rotation synthesises three ATP molecules.}
  \label{fig:atpsynthase}
\end{figure}

The F$_0$ motor is driven by proton flow through the $a$-$c$ interface.
Each proton captured by a $c$-subunit carboxylate (Asp$_{61}$ in
\Ecoli) causes the $c$-ring to advance by one subunit. A complete
revolution requires as many protons as there are $c$-subunits.
The $c$-ring stoichiometry varies: 8 in bovine and human mitochondria,
10 in yeast, 14--15 in plant chloroplasts.

\subsection{ATP synthesis stoichiometry}

Each full $\gamma$ rotation synthesises 3~ATP (one per $\beta$ subunit
per 360$^\circ$). With $c = 8$~subunits (mammalian):
\[
  \text{H}^+/\text{ATP} = \frac{c}{3} = \frac{8}{3} \approx 2.67.
\]
One additional proton is consumed per ATP for phosphate import
(the phosphate carrier co-transports H$^+$), giving an effective
$\text{H}^+/\text{ATP} \approx 3.67$.

Using 10~H$^+$ from NADH oxidation:
\[
  \text{ATP per NADH} = \frac{10}{3.67} \approx 2.73.
\]
For FADH$_2$ (6~H$^+$): ATP per FADH$_2 \approx 1.64$.

\textbf{P/O ratios} (ATP per oxygen atom reduced):
\begin{itemize}
  \item NADH: $\approx 2.5$~ATP (traditional textbook value: 3)
  \item FADH$_2$: $\approx 1.5$~ATP (traditional: 2)
  \item Glucose (complete oxidation): $\approx 30$--32 ATP (traditional: 36--38)
\end{itemize}
The discrepancy from older values arises because the traditional
values assumed an integer P/O ratio and ignored the proton cost of
the phosphate carrier.

% ============================================================
\section{Metabolic flux control}
% ============================================================

\subsection{Flux control coefficients}

Metabolic control analysis (MCA), developed by Kacser and Burns (1973)
and Heinrich and Rapoport (1974), provides a rigorous framework for
asking: how does a perturbation to one enzyme activity affect the
steady-state flux through a pathway?

The \textbf{flux control coefficient} of enzyme $i$ on pathway flux $J$:
\begin{equation}
  C_i^J = \frac{\partial J/J}{\partial E_i/E_i}
        = \frac{\partial \ln J}{\partial \ln E_i},
  \label{eq:fcc}
\end{equation}
the fractional change in flux per fractional change in enzyme activity.

The \textbf{summation theorem}:
\begin{equation}
  \sum_i C_i^J = 1.
  \label{eq:summation}
\end{equation}
This means flux control is shared among all enzymes in the pathway.
No single enzyme must carry $C_i^J = 1$; in practice, most enzymes
in a linear pathway have $C_i^J \ll 1$ and control is distributed.

\subsection{Why the rate-limiting step is a simplification}

The classical notion of a single rate-limiting step corresponds to the
special case where one enzyme has $C_i^J \approx 1$ and all others
$\approx 0$. Experimentally, this is rare. In glycolysis, for example,
phosphofructokinase-1 (PFK-1) has the largest control coefficient
($C \approx 0.3$--$0.6$ depending on conditions), but other steps
share the remaining control. Statements like ``hexokinase limits
glycolysis'' are context-dependent approximations.

% ============================================================
\section{Signal transduction}
% ============================================================

\subsection{Receptor-ligand binding}

For a simple receptor $R$ binding ligand $L$:
\[
  \mathrm{R} + \mathrm{L} \rightleftharpoons \mathrm{RL},
  \qquad \Kd = \frac{[\mathrm{R}][\mathrm{L}]}{[\mathrm{RL}]}.
\]
Fractional occupancy:
\[
  \theta = \frac{[\mathrm{L}]}{[\mathrm{L}] + \Kd}.
\]
The dose-response relation is hyperbolic with EC$_{50} = \Kd$ for
simple 1:1 binding. For systems where receptor occupancy must
cross a threshold to trigger a response, the steepness of the
response matters as much as the EC$_{50}$.

\subsection{Signal amplification}

Each step in a signaling cascade can amplify the signal. An RTK
phosphorylated on $n$ tyrosines can recruit $n$ adaptor molecules;
each adaptor can activate $m$ Ras molecules; each active Ras activates
one Raf kinase that phosphorylates many MEK molecules; each MEK
phosphorylates many ERK molecules. The cascade provides signal
amplification and noise filtering.

\begin{figure}[H]
  \centering
  \includegraphics[width=0.60\textwidth]{mapk_signaling_cascade_1786956266628.jpg}
  \caption{The RTK-Ras-MAPK signaling cascade. EGF binding induces
           RTK dimerisation and autophosphorylation. The adaptor complex
           Grb2-SOS activates Ras (GDP $\to$ GTP). Ras-GTP activates
           the three-kinase Raf-MEK-ERK module. ERK enters the nucleus
           and phosphorylates transcription factors. Negative feedback
           from ERK to SOS attenuates the signal.}
  \label{fig:mapk}
\end{figure}

\subsection{Ultrasensitivity and switch-like responses}

Many signaling responses are switch-like: they are OFF below a
threshold and ON above it. This arises from:
\begin{enumerate}
  \item \textbf{Cooperativity} (Hill coefficient $n > 1$): as derived
        in Module~1 for the Hill equation.
  \item \textbf{Zeroth-order ultrasensitivity} (Goldbeter-Koshland,
        1981): if kinase and phosphatase both operate near saturation
        ($[\text{substrate}] \gg K_M$), small changes in kinase/phosphatase
        ratio produce large all-or-none changes in phosphorylation.
  \item \textbf{Dual phosphorylation}: two sequential phosphorylation
        events required for activation (e.g., MEK requires
        phospho-Thr and phospho-Ser) generate ultrasensitivity
        equivalent to $n \approx 2$.
\end{enumerate}

The \textbf{Goldbeter-Koshland function} for a substrate phosphorylated
by a kinase $V_1$ and dephosphorylated by a phosphatase $V_2$, both
operating near saturation:
\begin{equation}
  \mathrm{W}^\ast = \frac{V_1/V_2 - 1 + \sqrt{(V_1/V_2 - 1)^2 + 4(V_1/V_2)K_2/K_1}}
                         {2(1 - K_2/K_1)},
  \label{eq:gk}
\end{equation}
where $W^\ast$ is the fraction phosphorylated, $K_1 = K_M^\text{kinase}$,
$K_2 = K_M^\text{phosphatase}$. When both $K_1, K_2 \to 0$
(zeroth-order regime), $W^\ast$ switches from 0 to 1 at $V_1 = V_2$
with infinite steepness.

% ============================================================
\section{Code lab: enzyme kinetics simulation}
% ============================================================

\begin{codelab}{Michaelis-Menten kinetics with competitive, uncompetitive, and mixed inhibition}
  Computes $v$ vs.\ $[\mathrm{S}]$ and the Lineweaver-Burk transformation
  for three inhibitor types, and fits simulated data to recover $K_M$
  and $V_{\max}$.
\end{codelab}

\begin{lstlisting}[language=Python,
  caption={Enzyme kinetics simulation (Python 3, NumPy, Matplotlib).},
  label={lst:enzkinetics}]
import numpy as np
import matplotlib.pyplot as plt

# -------------------------------------------------------
# Enzyme parameters
# -------------------------------------------------------
Vmax = 100.0    # umol/min/mg
Km   = 0.5      # mM
I    = 1.0      # mM inhibitor concentration
Ki   = 0.3      # mM  -- inhibits free E
Kip  = 0.6      # mM  -- inhibits ES (uncompetitive / mixed component)

S = np.linspace(0.01, 10.0, 500)  # mM

def mm_rate(S, Vmax, Km):
    return Vmax * S / (S + Km)

def competitive(S, Vmax, Km, I, Ki):
    """K_M apparent increases; V_max unchanged."""
    Km_app = Km * (1 + I/Ki)
    return Vmax * S / (S + Km_app)

def uncompetitive(S, Vmax, Km, I, Kip):
    """Both K_M and V_max decrease by same factor."""
    alpha = 1 + I/Kip
    return (Vmax/alpha) * S / (S + Km/alpha)

def mixed(S, Vmax, Km, I, Ki, Kip):
    """General case: alpha affects K_M, alpha' affects V_max."""
    alpha  = 1 + I/Ki
    alphap = 1 + I/Kip
    return Vmax * S / (alphap * S + alpha * Km)

v_none  = mm_rate(S, Vmax, Km)
v_comp  = competitive(S, Vmax, Km, I, Ki)
v_uncomp = uncompetitive(S, Vmax, Km, I, Kip)
v_mixed = mixed(S, Vmax, Km, I, Ki, Kip)

# -------------------------------------------------------
# Plot 1: v vs [S]
# -------------------------------------------------------
fig, axes = plt.subplots(1, 2, figsize=(13, 5))

ax = axes[0]
ax.plot(S, v_none,   'k-',   lw=2.0, label='No inhibitor')
ax.plot(S, v_comp,   'b--',  lw=1.8, label='Competitive')
ax.plot(S, v_uncomp, 'r-.',  lw=1.8, label='Uncompetitive')
ax.plot(S, v_mixed,  'g:',   lw=2.0, label='Mixed')
ax.axhline(Vmax,     color='gray', ls=':', lw=1)
ax.set_xlabel('[S] (mM)')
ax.set_ylabel('v (umol/min/mg)')
ax.set_title('Michaelis-Menten with inhibition')
ax.legend()
ax.grid(True, alpha=0.3)

# -------------------------------------------------------
# Plot 2: Lineweaver-Burk (1/v vs 1/[S])
# -------------------------------------------------------
S_lb  = S[S > 0.1]   # avoid 1/S -> infinity
ax = axes[1]
for v, label, ls in [(v_none[S>0.1],  'No inhibitor', 'k-'),
                     (v_comp[S>0.1],   'Competitive',  'b--'),
                     (v_uncomp[S>0.1], 'Uncompetitive','r-.'),
                     (v_mixed[S>0.1],  'Mixed',        'g:')]:
    ax.plot(1/S_lb, 1/v, ls, lw=1.8, label=label)
ax.set_xlim(-4, 12)
ax.set_ylim(-0.01, 0.08)
ax.axhline(0, color='gray', lw=0.5)
ax.axvline(0, color='gray', lw=0.5)
ax.set_xlabel('1/[S] (mM$^{-1}$)')
ax.set_ylabel('1/v')
ax.set_title('Lineweaver-Burk plot')
ax.legend(fontsize=9)
ax.grid(True, alpha=0.3)

plt.tight_layout()
plt.savefig('enzyme_inhibition.pdf', dpi=150)
plt.show()

# -------------------------------------------------------
# Print apparent kinetic parameters for each inhibitor type
# -------------------------------------------------------
alpha   = 1 + I/Ki
alphap  = 1 + I/Kip
print(f"\nNo inhibitor:    Km = {Km:.2f} mM,  Vmax = {Vmax:.1f}")
print(f"Competitive:     Km = {Km*alpha:.2f} mM,  Vmax = {Vmax:.1f}")
print(f"Uncompetitive:   Km = {Km/alphap:.2f} mM,  Vmax = {Vmax/alphap:.1f}")
print(f"Mixed:           Km = {Km*alpha/alphap:.2f} mM,  Vmax = {Vmax/alphap:.1f}")
\end{lstlisting}

\textbf{Results.}
\begin{itemize}
  \item Competitive inhibition: $K_M^\text{app} = 0.5(1 + 1/0.3) = 2.17$~mM;
        $V_{\max}$ unchanged at 100.
  \item Uncompetitive: $K_M^\text{app} = 0.5/(1 + 1/0.6) = 0.18$~mM;
        $V_{\max}^\text{app} = 100/(1 + 1/0.6) = 37.5$.
  \item Mixed: $K_M^\text{app} = 0.5\alpha/\alpha' = 0.81$~mM;
        $V_{\max}^\text{app} = 100/\alpha' = 37.5$.
\end{itemize}

% ============================================================
%  PROBLEM SET
% ============================================================
\newpage
\section{Problem set: Module 2}

\noindent Attempt all problems. Show all working.

% ------------------------------------------------------------
\begin{problem}{\stars{2}\enspace\disciplinetag{Enzyme Kinetics}\enspace
  Michaelis-Menten from initial rate data}

  An enzyme at $[\mathrm{E}]_0 = 5$~nM gives the following
  initial velocities:

  \begin{center}
  \begin{tabular}{cc}
  \toprule
  $[\mathrm{S}]$ (mM) & $v$ ($\mu$M~min$^{-1}$) \\
  \midrule
  0.1  & 1.37 \\
  0.5  & 4.55 \\
  1.0  & 6.25 \\
  2.0  & 7.69 \\
  5.0  & 8.93 \\
  \bottomrule
  \end{tabular}
  \end{center}

  \begin{enumerate}
    \item Plot $1/v$ vs $1/[\mathrm{S}]$ and determine $V_{\max}$ and
          $K_M$ from the intercepts.
    \item Calculate $k_\text{cat}$ and the catalytic efficiency
          $k_\text{cat}/K_M$.
    \item Is this enzyme catalytically perfect? Justify with reference
          to the diffusion limit.
  \end{enumerate}
\end{problem}

\begin{solution}
  \textbf{(a)} Lineweaver-Burk. Invert the data:

  \begin{center}
  \begin{tabular}{cc}
  \toprule
  $1/[\mathrm{S}]$ (mM$^{-1}$) & $1/v$ ($\mu$M$^{-1}$~min) \\
  \midrule
  10.0 & 0.730 \\
   2.0 & 0.220 \\
   1.0 & 0.160 \\
   0.5 & 0.130 \\
   0.2 & 0.112 \\
  \bottomrule
  \end{tabular}
  \end{center}

  A linear fit gives $y$-intercept $= 1/V_{\max} = 0.100$, so
  $V_{\max} = 10.0~\mu$M~min$^{-1}$.
  Slope $= K_M/V_{\max} = 0.063$~mM~$\mu$M$^{-1}$~min, so
  $K_M = 0.063 \times 10.0 = 0.63$~mM.

  \textbf{(b)} $k_\text{cat} = V_{\max}/[\mathrm{E}]_0 =
  10.0~\mu\text{M~min}^{-1} / 5~\text{nM} =
  10^{-5}/5\times10^{-9}~\text{s}^{-1}/60 = 33.3$~s$^{-1}$.\\
  $k_\text{cat}/K_M = 33.3 / (0.63\times10^{-3}) = 5.3\times10^4$~M$^{-1}$s$^{-1}$.

  \textbf{(c)} No. The diffusion limit is $\sim10^8$--$10^9$~M$^{-1}$s$^{-1}$.
  This enzyme is approximately $10^4$-fold below the diffusion limit,
  so it is not catalytically perfect.
\end{solution}

% ------------------------------------------------------------
\begin{problem}{\stars{3}\enspace\disciplinetag{Inhibition}\enspace
  Identifying inhibitor type from kinetic data}

  An enzyme has $V_{\max} = 80~\mu$M~min$^{-1}$ and
  $K_M = 0.4$~mM. In the presence of inhibitor at $[\mathrm{I}] = 2$~mM:
  $V_{\max}^\text{app} = 80~\mu$M~min$^{-1}$, $K_M^\text{app} = 1.2$~mM.

  \begin{enumerate}
    \item Identify the type of inhibition.
    \item Calculate $K_i$.
    \item An alternative inhibitor at the same 2~mM gives
          $V_{\max}^\text{app} = 40$ and $K_M^\text{app} = 0.2$~mM.
          What type is this, and what is $K_i'$?
  \end{enumerate}
\end{problem}

\begin{solution}
  \textbf{(a)} $V_{\max}$ is unchanged but $K_M$ increases: this is
  \textbf{competitive} inhibition.

  \textbf{(b)} $K_M^\text{app} = K_M(1 + [\mathrm{I}]/K_i)$:
  \[
    1.2 = 0.4\left(1 + \frac{2}{K_i}\right)
    \implies 1 + 2/K_i = 3 \implies K_i = 1.0~\text{mM}.
  \]

  \textbf{(c)} Both $V_{\max}$ and $K_M$ decrease by the same factor
  $(80 \to 40$ and $0.4 \to 0.2$, both halved): this is
  \textbf{uncompetitive} inhibition.\\
  $\alpha' = 1 + [\mathrm{I}]/K_i' = V_{\max}/V_{\max}^\text{app} = 80/40 = 2$:
  \[
    K_i' = \frac{2.0}{2 - 1} = 2.0~\text{mM}.
  \]
\end{solution}

% ------------------------------------------------------------
\begin{problem}{\stars{3}\enspace\disciplinetag{Bioenergetics}\enspace
  In vivo ATP hydrolysis free energy}

  At 37$^\circ$C in a hepatocyte: $[\mathrm{ATP}] = 3.5$~mM,
  $[\mathrm{ADP}] = 1.5$~mM, $[\mathrm{P}_i] = 4.0$~mM.
  Use $\dG^{\circ'} = -30.5$~kJ~mol$^{-1}$,
  $R = 8.314$~J~mol$^{-1}$~K$^{-1}$.

  \begin{enumerate}
    \item Calculate $\dG$ for ATP hydrolysis under these conditions.
    \item How many $k_\mathrm{B}T$ does this correspond to? ($k_\mathrm{B}T
          = RT/N_A = RT$ per molecule $\approx 2.58$~kJ~mol$^{-1}$ at 37$^\circ$C.)
    \item The plasma membrane Na$^+$/K$^+$ ATPase moves 3~Na$^+$
          out and 2~K$^+$ in per ATP hydrolysed. Against a
          $\Delta\psi = -70$~mV membrane potential, estimate whether
          this is thermodynamically feasible.
  \end{enumerate}
\end{problem}

\begin{solution}
  \textbf{(a)}
  \[
    \dG = -30{,}500 + 8.314 \times 310
          \times \ln\!\left(\frac{1.5\times10^{-3}\times4.0\times10^{-3}}
                                  {3.5\times10^{-3}}\right).
  \]
  \[
    = -30{,}500 + 2577 \times \ln(1.71\times10^{-3})
    = -30{,}500 + 2577 \times (-6.37)
    = -30{,}500 - 16{,}410 \approx -46.9~\text{kJ~mol}^{-1}.
  \]

  \textbf{(b)} $-46.9 / 2.58 = 18.2~k_\mathrm{B}T$.

  \textbf{(c)} The free energy cost of moving one ion against a membrane
  potential $\Delta\psi$ is $\dG_\text{elec} = zF\Delta\psi$ per mole.
  Moving 3~Na$^+$ outward (against $\Delta\psi = -70$~mV, so moving
  positive charge out costs energy):
  \[
    \dG_\text{Na} = 3 \times 96{,}485 \times 0.070 = +20.3~\text{kJ~mol}^{-1}.
  \]
  Moving 2~K$^+$ inward (K$^+$ enters, negative membrane potential
  provides driving force, but the concentration gradient and
  electrochemical potential must be considered; approximate as
  $\approx +6$~kJ~mol$^{-1}$ net cost for 2~K$^+$).
  Total ion transport cost $\approx 26$~kJ~mol$^{-1}$.
  ATP hydrolysis provides $-46.9$~kJ~mol$^{-1} > 26$~kJ~mol$^{-1}$:
  thermodynamically feasible with $\sim 20$~kJ~mol$^{-1}$ to spare.
\end{solution}

% ------------------------------------------------------------
\begin{problem}{\stars{3}\enspace\disciplinetag{Oxidative Phosphorylation}\enspace
  Proton motive force calculation}

  Mitochondria in an actively respiring cell maintain $\Delta\psi =
  -175$~mV and $\Delta\mathrm{pH} = 0.8$ (matrix more alkaline).

  \begin{enumerate}
    \item Calculate the PMF at 37$^\circ$C.
    \item Calculate the free energy per proton translocated from IMS
          to matrix.
    \item Using a $c$-ring stoichiometry of 8 subunits and accounting
          for the phosphate carrier (1 additional H$^+$/ATP), calculate
          the maximum number of ATP synthesised per NADH oxidised.
  \end{enumerate}
\end{problem}

\begin{solution}
  \textbf{(a)}
  \[
    \Delta p = -175 - 61.5 \times 0.8 = -175 - 49.2 = -224.2~\text{mV}.
  \]

  \textbf{(b)}
  \[
    \dG_{\text{H}^+} = F \times \Delta p = 96{,}485 \times (-0.2242)
    = -21.6~\text{kJ~mol}^{-1}.
  \]

  \textbf{(c)} NADH oxidation via Complexes I, III, IV pumps 10~H$^+$.
  H$^+$ per ATP (c-ring): $8/3 = 2.67$. Plus 1 for phosphate carrier:
  $2.67 + 1 = 3.67$~H$^+$ per ATP.
  \[
    \text{ATP/NADH} = \frac{10}{3.67} = 2.72.
  \]
\end{solution}

% ------------------------------------------------------------
\begin{problem}{\stars{3}\enspace\disciplinetag{Bioenergetics}\enspace
  Thermodynamic coupling: active transport}

  The sodium-glucose cotransporter SGLT1 imports 2~Na$^+$ and
  1 glucose per cycle. The sodium electrochemical gradient provides
  the driving force.

  Given: intracellular $[\mathrm{Na}^+]_i = 12$~mM, extracellular
  $[\mathrm{Na}^+]_o = 145$~mM, $\Delta\psi = -70$~mV, $T = 310$~K.
  Extracellular glucose $= 5$~mM; intracellular glucose $= 0.1$~mM.

  \begin{enumerate}
    \item Calculate the free energy of importing 2~Na$^+$ from outside
          to inside (concentration term + electrical term).
    \item Calculate the free energy required to import 1 glucose
          against its concentration gradient.
    \item Is the cotransport thermodynamically feasible?
  \end{enumerate}
\end{problem}

\begin{solution}
  \textbf{(a)} Free energy per mole of Na$^+$ transported inward:
  \[
    \dG_{\text{Na}^+} = RT\ln\frac{[\text{Na}^+]_i}{[\text{Na}^+]_o}
                        + zF\Delta\psi.
  \]
  \[
    = 8.314\times310\times\ln\frac{12}{145}
      + 1\times96{,}485\times(-0.070)
  \]
  \[
    = 2577\times(-2.49) - 6{,}754 = -6{,}416 - 6{,}754 = -13{,}170~\text{J~mol}^{-1}
    = -13.2~\text{kJ~mol}^{-1}.
  \]
  For 2~Na$^+$: $\dG = 2 \times (-13.2) = -26.4$~kJ~mol$^{-1}$.

  \textbf{(b)} Glucose import against its gradient (uncharged, $z=0$):
  \[
    \dG_\text{glucose} = RT\ln\frac{0.1}{5} = 2577\times\ln(0.02)
    = 2577\times(-3.91) = -10{,}076~\text{J~mol}^{-1}.
  \]
  This is negative (glucose is being imported down its gradient from
  5~mM to 0.1~mM), so no energy is required — wait, glucose moves
  from outside ($5$~mM) to inside ($0.1$~mM): \textbf{downhill}.
  Let us reconsider: $\dG = RT\ln([\text{in}]/[\text{out}]) =
  RT\ln(0.1/5) = -10.1$~kJ~mol$^{-1}$. Glucose import is spontaneous.

  \textbf{(c)} Total $\dG = -26.4 + (-10.1) = -36.5$~kJ~mol$^{-1} < 0$.
  Both Na$^+$ and glucose move downhill electrochemically, so
  transport is spontaneous. (SGLT1 operates against glucose gradient
  in the absorptive epithelium; the assumption here of higher
  extracellular glucose shows the driving force is entirely from Na$^+$.)
\end{solution}

% ------------------------------------------------------------
\begin{problem}{\stars{4}\enspace\disciplinetag{Enzyme Kinetics}\enspace
  Enzyme efficiency and transition-state theory}

  The activation energy for the uncatalysed hydrolysis of
  phosphodiester bonds is $\dG^\ddagger_\text{uncat} = 130$~kJ~mol$^{-1}$.
  Ribonuclease A lowers this to $\dG^\ddagger_\text{cat} = 55$~kJ~mol$^{-1}$.

  \begin{enumerate}
    \item Calculate the rate enhancement
          $k_\text{cat}/k_\text{uncat}$ at 310~K using the
          Arrhenius approximation
          $k \propto \exp(-\dG^\ddagger/RT)$.
    \item Transition-state stabilisation $= \dG^\ddagger_\text{uncat}
          - \dG^\ddagger_\text{cat}$. Convert to $k_\mathrm{B}T$ units.
    \item The uncatalysed rate constant is $k_\text{uncat} = 6\times10^{-6}$~s$^{-1}$
          and $k_\text{cat} = 1\times10^3$~s$^{-1}$ for ribonuclease A.
          Calculate $k_\text{cat}/K_M$ given $K_M = 1\times10^{-3}$~M
          and compare with the diffusion limit.
  \end{enumerate}
\end{problem}

\begin{solution}
  \textbf{(a)}
  \[
    \frac{k_\text{cat}}{k_\text{uncat}}
    = \exp\!\left(\frac{\dG^\ddagger_\text{uncat} - \dG^\ddagger_\text{cat}}{RT}\right)
    = \exp\!\left(\frac{(130-55)\times10^3}{8.314\times310}\right)
    = e^{29.1} \approx 4.5\times10^{12}.
  \]

  \textbf{(b)} Stabilisation $= 75$~kJ~mol$^{-1}$.
  In $k_\mathrm{B}T$ units: $75{,}000/2577 = 29.1~k_\mathrm{B}T$.

  \textbf{(c)}
  \[
    k_\text{cat}/K_M = 10^3 / 10^{-3} = 10^6~\text{M}^{-1}\text{s}^{-1}.
  \]
  Below the diffusion limit ($10^8$--$10^9$~M$^{-1}$s$^{-1}$)
  by $\sim$100--1000-fold. Ribonuclease A is not catalytically perfect.
\end{solution}

% ------------------------------------------------------------
\begin{problem}{\stars{3}\enspace\disciplinetag{Metabolic Control}\enspace
  Flux control coefficients and the summation theorem}

  A linear pathway $A \to B \to C \to D$ has three enzymes $E_1$,
  $E_2$, $E_3$ with measured flux control coefficients
  $C_1^J = 0.6$ and $C_2^J = 0.3$.

  \begin{enumerate}
    \item What is $C_3^J$?
    \item Enzyme activity $E_1$ is increased by 50\% using a genetic
          approach. Predict the new flux, given initial $J_0 = 100$
          (arbitrary units). State any assumption.
    \item A drug inhibits $E_2$ by 20\%. Estimate the fractional
          change in flux.
    \item Why does the enzyme with the highest control coefficient
          not necessarily correspond to the enzyme with the lowest
          activity in the pathway?
  \end{enumerate}
\end{problem}

\begin{solution}
  \textbf{(a)} Summation theorem: $\sum C_i^J = 1$.
  \[
    C_3^J = 1 - 0.6 - 0.3 = 0.1.
  \]

  \textbf{(b)} By definition of $C_1^J = (\partial J/J)/(\partial E_1/E_1)$,
  for small perturbations:
  \[
    \frac{\delta J}{J} \approx C_1^J \cdot \frac{\delta E_1}{E_1}
    = 0.6 \times 0.5 = 0.30.
  \]
  New flux: $J \approx 100 \times 1.30 = 130$. Assumption: linear
  regime (perturbation is small relative to the system's nonlinearity).
  A 50\% increase may push beyond the linear regime.

  \textbf{(c)} Inhibiting $E_2$ by 20\%:
  \[
    \frac{\delta J}{J} = C_2^J \cdot \frac{\delta E_2}{E_2}
    = 0.3 \times (-0.20) = -0.06.
  \]
  Flux decreases by approximately 6\%.

  \textbf{(d)} Flux control coefficient depends on the kinetic
  parameters of the enzyme \emph{in context} (concentrations of
  substrates and products relative to $K_M$), not on the enzyme's
  absolute activity. An enzyme at saturation ($[\mathrm{S}] \gg K_M$)
  has near-zero control coefficient regardless of its $V_{\max}$
  because increasing its activity further does not change the rate
  through the saturated enzyme.
\end{solution}

% ------------------------------------------------------------
\begin{problem}{\stars{3}\enspace\disciplinetag{Signal Transduction}\enspace
  Dose-response and receptor occupancy}

  An RTK has $\Kd = 2$~nM for its ligand EGF. The downstream MAPK
  response requires $\geq 50\%$ receptor occupancy for activation.

  \begin{enumerate}
    \item What EGF concentration achieves 50\% receptor occupancy?
    \item A tumour mutation constitutively dimerises the RTK, reducing
          the effective $\Kd$ to 0.05~nM. At a physiological EGF
          concentration of 1~nM, what is the receptor occupancy for
          wild-type vs.\ mutant?
    \item The MAPK cascade has 3 sequential phosphorylation steps
          (Raf, MEK, ERK), each with Hill coefficient $n=1$ and
          EC$_{50} = 0.5$ of its upstream input. If Raf activation
          is $\theta_\text{Raf} = 0.5$, calculate ERK activation.
  \end{enumerate}
\end{problem}

\begin{solution}
  \textbf{(a)} For $\theta = 0.5$: $[\mathrm{EGF}] = \Kd = 2$~nM.

  \textbf{(b)} Wild-type ($\Kd = 2$~nM, $[\mathrm{EGF}] = 1$~nM):
  \[
    \theta_\text{WT} = \frac{1}{1+2} = 0.33.
  \]
  Mutant ($\Kd = 0.05$~nM):
  \[
    \theta_\text{mut} = \frac{1}{1 + 0.05} = 0.95.
  \]
  Wild-type receptor is 33\% occupied (below the 50\% threshold);
  mutant is 95\% occupied — constitutively active.

  \textbf{(c)} Each cascade step compresses the dynamic range.
  With EC$_{50} = 0.5$ of input signal and $n=1$:
  MEK input $= \theta_\text{Raf} = 0.5$, so
  $\theta_\text{MEK} = 0.5/(0.5 + 0.5) = 0.5$.
  ERK input $= 0.5$: $\theta_\text{ERK} = 0.5$.
  With $n=1$ and symmetric EC$_{50}$, the cascade does not amplify at
  this operating point. Amplification requires $n>1$ or asymmetric
  EC$_{50}$.
\end{solution}

% ------------------------------------------------------------
\begin{problem}{\stars{4}\enspace\disciplinetag{Oxidative Phosphorylation}
  \enspace\disciplinetag{Bioenergetics}\enspace
  ATP yield from glucose}

  Complete oxidation of one glucose molecule produces:
  10 NADH, 2 FADH$_2$, and 2 substrate-level ATP (from glycolysis
  and the TCA cycle). Use P/O ratios of 2.5 (NADH) and 1.5 (FADH$_2$).

  \begin{enumerate}
    \item Calculate the total ATP yield per glucose.
    \item The standard free energy of glucose oxidation:
          $\dG^{\circ'} = -2870$~kJ~mol$^{-1}$.
          Assuming $|\dG_\text{ATP}| = 50$~kJ~mol$^{-1}$ per ATP
          under cellular conditions, calculate the thermodynamic
          efficiency of oxidative phosphorylation.
    \item Under anaerobic conditions, only glycolysis operates
          (net 2 ATP per glucose). Calculate the ratio of aerobic
          to anaerobic ATP yield and the corresponding
          thermodynamic efficiency of fermentation, noting
          that $\dG^{\circ'}(\text{glucose} \to 2 \text{ lactate})
          = -196$~kJ~mol$^{-1}$.
  \end{enumerate}
\end{problem}

\begin{solution}
  \textbf{(a)}
  \[
    \text{ATP} = 10 \times 2.5 + 2 \times 1.5 + 2 = 25 + 3 + 2 = 30~\text{ATP}.
  \]

  \textbf{(b)} Free energy captured:
  $30 \times 50 = 1500$~kJ~mol$^{-1}$.\\
  Efficiency: $\eta = 1500/2870 = 52\%$.

  \textbf{(c)} Aerobic/anaerobic ratio: $30/2 = 15$.\\
  Fermentation efficiency: $(2 \times 50)/196 = 100/196 = 51\%$.
  Note: the efficiency appears comparable, but the \emph{power} output
  (ATP rate) of aerobic respiration is $\sim$15-fold higher.
\end{solution}

% ------------------------------------------------------------
\begin{problem}{\stars{5}\enspace\disciplinetag{Signal Transduction}
  \enspace\disciplinetag{Enzyme Kinetics}\enspace\disciplinetag{Bioenergetics}\enspace
  Oncogenic mutation: energetic and signaling consequences}

  A somatic mutation in Ras (G12V) locks it in the GTP-bound active
  state by impairing intrinsic GTPase activity. Wild-type Ras
  hydrolyses GTP with $k_\text{cat} = 0.02$~min$^{-1}$; the G12V
  mutant has $k_\text{cat} = 0.0002$~min$^{-1}$.

  \begin{enumerate}
    \item Calculate the mean lifetime of the active Ras-GTP state
          for wild-type and G12V at the single-molecule level.
    \item GAP (GTPase-activating protein) accelerates wild-type GTP
          hydrolysis by $10^5$-fold but is ineffective on G12V
          (mutation blocks the GAP-stabilised transition state).
          What is the wild-type $k_\text{cat}$ in the presence of GAP?
    \item Constitutively active Ras continuously drives the MAPK
          cascade. If ERK phosphorylation requires a sustained
          stimulus to activate transcription and each ERK-P molecule
          has a mean lifetime of $\tau = 3$~min (phosphatase action),
          and if 500 ERK molecules per cell must be phosphorylated
          simultaneously for threshold activation, estimate the
          minimum Ras-GTP concentration required given
          $k_\text{ERK,phos} = 0.1$~min$^{-1}$~molecule$^{-1}$~Ras.
    \item GTP hydrolysis consumes one GTP per Ras activation cycle.
          At a cytoplasmic $[\text{GTP}] = 0.5$~mM and
          $\dG_\text{GTP} = -51$~kJ~mol$^{-1}$, and assuming
          $10^6$ Ras molecules per cell actively cycling at $v = k_\text{cat,GAP}$,
          estimate the power dissipated by Ras-GAP catalysed GTP
          hydrolysis per cell (in fW).
  \end{enumerate}
\end{problem}

\begin{solution}
  \textbf{(a)} Mean lifetime = $1/k_\text{cat}$ (for first-order decay).\\
  Wild-type: $\tau_\text{WT} = 1/0.02~\text{min}^{-1} = 50$~min.\\
  G12V: $\tau_{G12V} = 1/0.0002 = 5000$~min $\approx$ 3.5~days.
  G12V Ras is effectively constitutively active on a cellular timescale.

  \textbf{(b)} GAP-stimulated $k_\text{cat}$:
  $0.02 \times 10^5 = 2000$~min$^{-1} \approx 33$~s$^{-1}$.

  \textbf{(c)} At steady state, the phosphorylation rate equals the
  dephosphorylation rate. If [Ras-GTP] molecules drive ERK phosphorylation
  at $k_\text{phos}$ per Ras:
  \[
    \frac{d[\text{ERK-P}]}{dt} = k_\text{phos} \times [\text{Ras-GTP}]
    - \frac{[\text{ERK-P}]}{\tau} = 0
  \]
  at steady state. Setting $[\text{ERK-P}] = 500$ molecules:
  \[
    [\text{Ras-GTP}] = \frac{[\text{ERK-P}]}{k_\text{phos}\times\tau}
    = \frac{500}{0.1 \times 3} = 1667~\text{molecules/cell}.
  \]

  \textbf{(d)} Hydrolysis rate per cell:
  \[
    \text{Rate} = 10^6 \times k_\text{cat,GAP} = 10^6 \times 2000~\text{min}^{-1}
    = 2\times10^9~\text{GTP~min}^{-1} = 3.33\times10^7~\text{s}^{-1}.
  \]
  Convert to moles per second: $3.33\times10^7 / 6.022\times10^{23}
  = 5.5\times10^{-17}$~mol~s$^{-1}$.\\
  Power $= \dot{n} \times |\dG| = 5.5\times10^{-17} \times 51{,}000
  = 2.8\times10^{-12}$~W $= 2.8$~pW per cell.
  This is a significant fraction of the estimated total cell power
  ($\sim1$--10~pW for a resting cell).
\end{solution}

% ============================================================
%  END OF MODULE 2
% ============================================================

\end{document}
